\chapter{Adding Thunks to CakeML}
\label{chap:cakeml}

Moving downwards on the compilation pipeline, this chapter describes our work
for the last part of the thesis, adding thunks to CakeML. The CakeML compiler
was not built for a lazy language, so it rightly doesn't consider thunks
natively. However, when considered as a target for PureCake, it is critical that
it can handle thunks effectively, since they are the main source of inefficiency
between lazy and strict programs.

\section{Taking Advantage of Thunks in the Runtime}

CakeML, like every high-level functional language, has a garbage collected
runtime. Each time a certain threshold of allocated memory is surpassed, the
program stops running and the garbage collector traverses the memory starting
from known used locations called roots, copying the reachable data into a new
heap and disposing of the old one. The details of garbage collection in CakeML
are described in \cite{sandberg2019verified}, but this mental model is enough
for our description here.

By having thunks represented explicitly in the CakeML source language, we can
make them explicit through all of CakeML's intermediate representations down to
the runtime. This would give the ability to the garbage collector to
discriminate between chunks of memory that represent thunks, which would allow
an important optimisation called updating \cite{seward1992generational}.

\begin{figure}[h]
  \centering
  \begin{tikzpicture}
    \tikzstyle{box} = [rectangle, text centered, draw=black]
    \node (box0) [box, draw=none] {$x$ : };
    \node (box1) [box, right=0cm of box0] { $ptr$ };
    \node (box2) [box, right=1cm of box1, draw=none] { $ptr$ : };
    \node (box3) [box, right=0cm of box2] { \lit{Evaluated}, 5 };
  \end{tikzpicture}
  \caption{Evaluated Thunk Before Optimisation}
  \label{fig:thunk}
\end{figure}

To illustrate the idea of the optimisation, suppose we have a thunk that has
evaluated to the value 5 stored in variable location $x$. The memory block at
$x$ will hold a pointer to the thunk block, which will hold the actual thunk
value. Since we now know that the two cells at $ptr$ are a thunk and not just an
array, we can instruct the garbage collector to replace the thunk with the
evaluated value when it encounters it. After the garbage collection, the block
at the $ptr$ location should be completely dropped, and the block in location
$x$ should just hold the value 5 directly.

This optimisation is crucial for both the memory and time efficiency of the
program. When an evaluated thunk is in memory two words of memory are wasted,
the one holding the pointer to the thunk and the \lit{Evaluated} tag. Getting
rid of these makes for more compact heaps that help the runtime efficiency of
the program. Furthermore, having the value directly available results in one
less indirection that the execution of the program has to traverse each time the
value is demanded, making accesses to it have the same performance
characteristics as if it were unboxed to begin with.

\todo[inline, size=normalsize] {
  More technical details will be inserted once all work has been completed on
  the implementation and proofs.
}
